% ===================================================================
%  Dodatek K2: Chiralna asymetria mas w TGP
%  TGP v1 — Sesja v41 (2026-03-30)
% ===================================================================
\section{Chiralna asymetria mas: $m_R \neq m_L$ z~substratu}%
\label{app:K2-chiralnosc}

Niniejszy dodatek formalizuje mechanizm \textbf{chiralnej asymetrii
mas} wynikaj\k{a}cy z~asymetrii potencja\l{}u $V(g)$ i~sprzę\.zenia
kinetycznego $K_{\rm sub}(g)$ solitonu TGP.
Wynik: stosunek mas $m_R/m_L = (A_{\rm tail}^R/A_{\rm tail}^L)^4$
jest nietrywialny, a~\textbf{bariera duchowa} przy $g^*$ naturalnie
\l{}amie chiralnosc bez pola Higgsa.

Wyniki opieraj\k{a} si\k{e} na skrypcie \texttt{ex108} (9/9 test\'ow
pozytywnych).

\medskip
\noindent\fbox{\parbox{\linewidth}{%
\textbf{Nota kanoniczna (2026-04-07).}
Niniejszy dodatek u\.zywa Formulacji~B ($g_0^e = 1{,}2492$, punkt duchowy $g^*$).
W~sformu\l{}owaniu kanonicznym: $K(g) = g^4$, $P(g) = (\beta/7)g^7 - (\gamma/8)g^8$,
$g_0^e = 0{,}86941$, brak punktu duchowego.
Algebra Koide'a i~r\'ownania master F1--F12 s\k{a} niezale\.zne od tego wyboru.}}
\medskip

% -------------------------------------------------------------------
\subsection{Asymetria potencja\l{}u $V(g)$}\label{app:K2-Vasym}
% -------------------------------------------------------------------

\begin{proposition}[Kubiczna asymetria potencja\l{}u]\label{prop:K2-Vasym}
Potencja\l{} solitonowy TGP:
\begin{equation}\label{eq:K2-V}
  V(g) = \frac{g^3}{3} - \frac{g^4}{4},
  \qquad
  V'(g) = g^2(1 - g),
  \qquad
  V(1) = \frac{1}{12},
\end{equation}
nie jest symetryczny wzgl\k{e}dem $g = 1$:
\begin{equation}\label{eq:K2-Vasym-expansion}
  V(1+\delta) - V(1-\delta)
  = \frac{V'''(1)}{3}\,\delta^3 + \mathcal{O}(\delta^5),
  \qquad
  V'''(1) = 2 - 6 = -4 \neq 0.
\end{equation}
Asymetria jest \textbf{kubiczna} ($\propto \delta^3$),
bo $V(g)$ ma oba cz\l{}ony $g^3$ i~$g^4$.
\end{proposition}

\begin{proof}
Bezpo\'srednia kalkulacja: $V''(g) = 2g - 3g^2$, $V'''(g) = 2 - 6g$,
$V'''(1) = -4$. Wyraz $V''(1) = -1$ (symetryczny) nie wchodzi
do r\'o\.znicy $V(1{+}\delta) - V(1{-}\delta)$.
\end{proof}

% -------------------------------------------------------------------
\subsection{Pary chiralne: kink-L i~antikink-R}\label{app:K2-pary}
% -------------------------------------------------------------------

\begin{definition}[Para chiralna]\label{def:K2-chiral-pair}
Dla danego parametru strzału $g_0$ definiujemy \textbf{parę chiralną}:
\begin{align}
  \text{Antikink-R (praworęczny):} &\quad g_0^R = 1 + \Delta,
    \qquad g(0) = g_0^R > 1,\quad g(\infty) = 1, \label{eq:K2-right}\\
  \text{Kink-L (leworęczny):} &\quad g_0^L = 1 - \Delta,
    \qquad g(0) = g_0^L < 1,\quad g(\infty) = 1. \label{eq:K2-left}
\end{align}
Oba profile spe\l{}niaj\k{a} to samo r\'ownanie ODE
z~pe\l{}nym sprzę\.zeniem substratowym $K_{\rm sub}(g) = g^2$:
\begin{equation}\label{eq:K2-ODE}
  K_{\rm sub}(g)\,g''
  + \frac{K_{\rm sub}'(g)}{2}\,(g')^2
  + \frac{2}{r}\,K_{\rm sub}(g)\,g'
  = V'(g).
\end{equation}
\end{definition}

\begin{proposition}[Chiralna asymetria ogon\'ow]\label{prop:K2-Atail-asym}
Dla pary chiralnej $(g_0^R, g_0^L) = (1+\Delta, 1-\Delta)$:
\begin{equation}\label{eq:K2-Atail-ratio}
  A_{\rm tail}^R(\Delta) \neq A_{\rm tail}^L(\Delta),
\end{equation}
a~stosunek chiralny:
\begin{equation}\label{eq:K2-rchiral}
  r_{\rm chiral}(\Delta)
  = \left(\frac{A_{\rm tail}^R}{A_{\rm tail}^L}\right)^{\!4}
  = \frac{m_R}{m_L},
\end{equation}
jest \textbf{monotonicznie rosn\k{a}cy} z~$\Delta > 0$
i~$r_{\rm chiral}(0) = 1$ (symetria przywr\'ocona).
Dla $\Delta = g_0^* - 1 \approx 0{,}249$ (elektron):
$r_{\rm chiral} \approx 2{,}03$.
\end{proposition}

\begin{proof}[Dow\'od (numeryczny)]
Skrypt \texttt{ex108} rozwi\k{a}zuje ODE~\eqref{eq:K2-ODE}
dla 5~warto\'sci $\Delta \in [0{,}05;\, 0{,}25]$ i~potwierdza:
(i)~$A_{\rm tail}^R > A_{\rm tail}^L$ dla ka\.zdego $\Delta > 0$;
(ii)~$r_{\rm chiral}$ monotonicznie rosn\k{a}cy;
(iii)~$r_{\rm chiral}(0{,}249) = 2{,}033$ (test T4, 9/9 PASS).
\'Zr\'od\l{}em asymetrii s\k{a} dwa czynniki dzia\l{}aj\k{a}ce
w~tym samym kierunku:
$V'''(1) = -4 < 0$ (potencja\l{}) i~$K_{\rm sub}(1+\delta)
= (1+\delta)^2 > (1-\delta)^2 = K_{\rm sub}(1-\delta)$ (kinetyczny).
\end{proof}

% -------------------------------------------------------------------
\subsection{Bariera duchowa i~\l{}amanie chiralno\'sci}\label{app:K2-ghost}
% -------------------------------------------------------------------

\begin{theorem}[Mechanizm chiralnej selekcji]\label{thm:K2-chiral-selection}
W~formalizacji $f(g) = 1 + 2\alpha\ln g$ (sek.~\ref{ssec:ghost-problem}):
\begin{enumerate}[nosep]
  \item Punkt duchowy: $g^* = \exp(-1/(2\alpha))
        \approx 0{,}779$ (dla $\alpha = 2$).
  \item Leworęczny kink elektronu: $g_0^L = 1 - (g_0^* - 1)
        = 2 - g_0^* \approx 0{,}751$.
  \item \textbf{$g_0^L < g^*$} --- leworęczny partner elektronu
        znajduje si\k{e} \textbf{za barier\k{a} duchow\k{a}}.
\end{enumerate}
W~opisie pe\l{}nym ($K_{\rm sub}(g) = g^2$) profil istnieje,
ale jego dynamika jest jako\'sciowo inna ni\.z prawor\k{e}cznego:
bariera $f(g^*) = 0$ dzia\l{}a jako \textbf{selektywna \'sciana},
przepuszczaj\k{a}c antikink-R (kt\'ory nie przekracza $g = 1$)
i~blokuj\k{a}c g\l{}\k{e}boki kink-L.
\end{theorem}

\begin{proof}
Bezpo\'srednie podstawienie:
$g_0^e = g_0^* = 1{,}2492$ (z~$\varphi$-FP, tw.~\ref{thm:J2-FP}),
$g_0^L = 2 - g_0^e = 0{,}7508$;
$g^* = e^{-1/4} = 0{,}7788$;
$g_0^L = 0{,}751 < g^* = 0{,}779$. \qed
\end{proof}

\begin{remark}[Chiralno\'s\'c bez pola Higgsa]\label{rem:K2-no-Higgs}
Mechanizm \l{}amania chiralno\'sci w~TGP nie wymaga pola Higgsa
ani \.zadnego dodatkowego pola ska\-lar\-ne\-go.
Wynika z~tych samych element\'ow, kt\'ore generuj\k{a} masy
(potencja\l{} $V(g)$, sprzę\.zenie $K_{\rm sub}(g)$)
i~jest fizyczn\k{a} konsekwencj\k{a} nieliniowej natury solitonu.
Jest to fundamentalna r\'o\.znica od Modelu Standardowego,
w~kt\'orym chiralno\'s\'c jest za\l{}o\.zona (leworęczne dublety,
praworęczne singlety), nie wynikowa.
\end{remark}

\begin{remark}[Falsyfikowalno\'s\'c]\label{rem:K2-falsyfikacja}
Predykcja $r_{\rm chiral} \approx 2{,}0$ mo\.ze by\'c
zweryfikowana przez:
\begin{itemize}[nosep]
  \item Pomiar prawor\k{e}cznych neutrin (je\'sli istniej\k{a}):
        ich masa powinna by\'c $\sim 2 \cdot m_\nu^L$.
  \item Asymetria p\k{e}d\'ow w~rozpadach $\beta$:
        efekt chiralny $\propto m_R/m_L - 1 \sim 1$.
\end{itemize}
Status: \textbf{Propozycja} (warunkowo na numeryce \texttt{ex108}).
\end{remark}
